\section{From a target position to random execution}
\label{sec:randomexecution}

A forecast of the reference price is not automatically the forecast
relevant to an order. A position target may fill partially, and the
conditional return of filled orders may differ from the unconditional
return. The earlier convex and causal results remain valid for their
stated executable-position objectives. This section makes a separate,
restricted random-execution cell explicit and identifies the moments a
forecast-to-order conversion actually requires.

\subsection{What changes relative to a fully filled target}
\citet[Section~37.1]{NoguerAlonso2026Execution} lists six modelling requirements---fill
law, impact and liquidity, information, competition, objective and
evaluation---and their status for a parent mandate. For a forecast-driven order
two of them change character. First, the relevant forecast is no longer the
conditional mean of the reference return but the fill-weighted moment
$\mathbb E[\Phi R]$, together with $\mathbb E\Phi$ and $\mathbb E\Phi^2$;
Section~\ref{sec:randomfill-rule} derives them. Second, position and capital
limits must hold for every possible executed position, not its mean
(Section~\ref{sec:fillbox}). The remaining four requirements carry over unchanged; in
particular the exogenous book law of Section~5 is a conditional modelling
assumption, not a model of how opponents react to the order policy.

\subsection{An exact decision rule with random fills}\label{sec:randomfill-rule}

Condition throughout on information available at the order decision.
There is an initial holding \(b\), an attempted increment \(u\) in a
closed interval \([\underline u,\overline u]\) containing zero, a fill
fraction \(\Phi\in[0,1]\), and a return \(R\). Execution occurs at the
start of this one-period cell at the reference price, with per-executed-unit
charge \(c\ge0\). The resulting position is \(b+\Phi u\).
A quadratic order charge \(\kappa u^2/2\), \(\kappa\ge0\), can be
included separately. A rebate alone is not to be substituted for a
nonnegative total \(c\) if that would violate this convexity assumption.

The joint law of \((\Phi,R)\) is held fixed for every admissible \(u\) in
this cell. This is a strong local assumption: changing direction, size,
price or venue generally changes it. One can restrict the interval to a
single side, or evaluate separate cells for different order instructions.
If fills occur at different prices or times, replace \(\Phi R\) by the
matched sum of executed quantities times their subsequent price changes
per unit attempted order, and rederive the inventory-risk term.

Define the conditional moments
\[
 p=\E\Phi,\qquad s_2=\E\Phi^2,\qquad m_F=\E[\Phi R],
 \qquad H=a s_2+\kappa>0,
\]
where \(a>0\) is the quadratic position charge. Relative to retaining
\(b\), expected risk-penalized gain is
\begin{equation}\label{eq:randomfillgain}
 G(u)=u(m_F-abp)-cp|u|-\tfrac12 H u^2.
\end{equation}
The target signal is \(m_F\), not \(p\E R\), and the risk coefficient
uses \(s_2\), not \(p^2\).

\begin{proposition}[A fill-weighted order rule]
\label{prop:fillweightedorder}
Under the preceding assumptions, let \(\mathcal S_d(x)=
\operatorname{sgn}(x)(|x|-d)_+\) and let
\(\Pi_{[\underline u,\overline u]}\) denote scalar clipping. The unique
optimal attempted increment is
\begin{equation}\label{eq:fillweightedorder}
 u^*=\Pi_{[\underline u,\overline u]}
       \left(\frac{\mathcal S_{cp}(m_F-abp)}{H}\right).
\end{equation}
If \(p,s_2\) are known but \(m_F\in[\widehat m_F-\varepsilon,
\widehat m_F+\varepsilon]\), the minimax gain relative to retaining
\(b\) is maximized by
\begin{equation}\label{eq:robustfillweightedorder}
 u^{\rm rob}=\Pi_{[\underline u,\overline u]}
       \left(\frac{\mathcal S_{cp+\varepsilon}(\widehat m_F-abp)}{H}\right).
\end{equation}
Its worst-case gain is nonnegative, since \(u=0\) is feasible.
\end{proposition}
\begin{proof}
Subtract \(bR-ab^2/2\) from
\((b+\Phi u)R-a(b+\Phi u)^2/2-c\Phi|u|-\kappa u^2/2\), and take
conditional expectations to obtain \eqref{eq:randomfillgain}. The
strictly concave scalar objective has the soft-threshold optimizer
in \eqref{eq:fillweightedorder}; clipping gives its maximum on the
interval. Minimizing its linear \(u m_F\) term over the mean interval
subtracts \(\varepsilon|u|\), giving
\eqref{eq:robustfillweightedorder}. Strict concavity and feasibility of
zero establish uniqueness and nonnegative worst-case gain.
\end{proof}

The elementary identity
\begin{equation}\label{eq:fillforecastbias}
 m_F=p\E R+\operatorname{Cov}(\Phi,R)
\end{equation}
shows how adverse selection enters the signal. If \(p>0\), then
\(m_F/p\) is the fill-weighted expected return, equal to a return
conditional on filling when \(\Phi\) is an indicator. Even under
independence, replacing a random position by its mean understates its
quadratic charge by
\(a u^2\operatorname{Var}(\Phi)/2\). This order rule is a specialization
of the earlier mixed-friction optimization; its contribution here is to
identify which execution moments must enter it.

For vector attempted orders and a diagonal random fill matrix \(D\),
with fixed joint law and symmetric risk matrix \(A\), the corresponding
quadratic coefficient is
\[
 H=\E[D^{\mathsf T}AD]+K,
 \qquad \eta=\E[D^{\mathsf T}(R-Ab)],
\]
where \(K\) is the quadratic order-charge matrix. Proportional charges
use each component's expected fill fraction. Off-diagonal entries retain
joint fills as well as risk correlations. The constrained convex problem
is available when \(H\) is positive definite; componentwise scalar
thresholding generally fails when it is not diagonal. This formula does
not model cross-impact or opponent response unless those terms are
explicitly supplied in the objective and law.

\subsection{A profitable price forecast can imply no buy order}\label{sec:couplingexample}

In a synthetic one-sided cell, let \(b=0\), \(u\in[0,1]\), \(a=0.10\),
\(c=0.01\), \(\kappa=0\), and let \(R\) equal \(0.14\) or \(-0.06\)
with equal probabilities. Thus \(\E R=0.04\). A Bernoulli fill has
\(p=s_2=0.5\). Under independence \(m_F=0.02\), so the rule attempts
\(u=0.30\) and predicts a gain of \(0.00225\).

Now couple every fill to the negative-return event. The return forecast,
fill probability and quadratic coefficient are unchanged, but
\(m_F=-0.03\). Sending the same \(0.30\) buy order produces expected
risk-penalized gain \(-0.01275\); the optimal admissible action is zero.
For a mixture with adverse-coupling weight \(\theta\),
\(m_F=0.02-0.05\theta\) and the optimal attempt is
\(u^*=(0.30-\theta)_+\). If only
\(0\le\theta\le\overline\theta\) is known, the maximin buy attempt is
\((0.30-\overline\theta)_+\). This last statement uses a one-sided
uncertainty interval and one-sided action set; it is not an assumption
that buy and sell fills have the same law.

\begin{table}[htbp]
\centering\small
\caption{Random-execution stress with identical marginal return and fill
laws. The last column evaluates the original \(u=0.30\) attempt under
each law. Values are synthetic objective units.}
\begin{tabular}{lrrr}
\toprule
Law & \(m_F\) & Optimal buy \(u^*\) & Gain at \(u=0.30\)\\
\midrule
\inputtable{tables/realism_rows.tex}
\bottomrule
\end{tabular}
\end{table}

The script \texttt{verification/realism\_stress.py} verifies the exact
example and compares the scalar rules with independent bounded
optimization for random discrete joint laws, position intervals and
uncertainty widths. A joint-moment confidence set must be calibrated;
one cannot relabel an unconditional mean interval as an interval for
\(m_F\). If \(p,s_2\) are uncertain too, they must join that set and the
simple formula requires a new robust optimization.

\subsection{An evaluation target that matches the order}

Compare unconditional-return prediction, joint fill/return prediction
and the induced feasible orders on the same chronological splits.
Report calibration of \(p,s_2,m_F\), actual completed positions, net cash
P\&L and the declared risk-penalized objective separately. Match each
return label to the execution price and horizon; do not use later
fills or quotes when constructing the earlier decision. Test rank,
latency, volatility, fees and response changes jointly, because separate
marginal accuracy checks cannot detect \eqref{eq:fillforecastbias}.
A target-position ablation is informative only when its full-execution
assumption is reported.

Finally, the sequential guarantee of Section~\ref{sec:certifiedposition}
uses a specified conditional exponential-moment bound. A random-filled
position changes its noise term. The earlier return-only variance proxy
cannot be carried over without proving the corresponding bound for joint
fill/return increments. Coverage of conditional decision gains,
control of realized noise, and a market counterfactual are three distinct
requirements.
